\begin{table}[H]
\centering
\caption{\textbf{Graph vs Subgraph Compression (C2: Interpretability-Oriented Compression).} Replacement measures fraction of end-to-end influence routed via features; Completeness measures fraction of incoming edge influence explained by features/tokens. Values are descriptive (no significance tests). \textsc{Pre-specified:} Metric definitions (Neuronpedia standard), subgraph evaluation procedure (influence-threshold selection via cumulative curve). \textsc{Exploratory:} Specific prompt set (selected to illustrate method across diverse domains). \emph{Note:} Completeness remains high (mean 0.83 vs 0.90 full graph), while Replacement drops (mean 0.54 vs 0.70), reflecting intentional compression that prioritizes legibility over exhaustive feature coverage.}
\label{tab:graph-subgraph}
\begin{tabular}{lrrrr}
\toprule
\textbf{Prompt} & \textbf{Graph Repl.} & \textbf{Subgraph Repl.} & \textbf{Graph Comp.} & \textbf{Subgraph Comp.} \\
\midrule
Austin (Dallas prompt)      & 0.72 & 0.57 & 0.90 & 0.83 \\
Oakland (entity swap)       & 0.70 & 0.57 & 0.90 & 0.83 \\
Michael Jordan (sports)     & 0.69 & 0.49 & 0.90 & 0.81 \\
Small opposite (general)    & 0.73 & 0.62 & 0.91 & 0.86 \\
Muscle diaphragm (anatomy)  & 0.63 & 0.45 & 0.87 & 0.79 \\
\midrule
Mean                        & 0.70 & 0.54 & 0.90 & 0.83 \\
\bottomrule
\end{tabular}
\end{table}


